\section{Cubic equations}

As soon as the genus is greater than 2, there are also cubic equations defining the Kummer va\-riety. To focus on cubic relations, we first note that the spaces $S^3 B_1$ and $B_3^+$ decompose under the action of the Heisenberg group into $2^g$ irreducible representations that are all isomorphic to~$B_1$ (see~\cite{vG} or \cite{Bea}); we are interested in studying the kernel of the map $ \Psi_3\colon S^3B_1\to B_3^+$. A~simple way of constructing cubic relations is considering quadratic relations. Quadratic relations exist whenever we have the vanishing of some even theta constants i.e., $ \theta\left[\begin{smallmatrix}\ep \\ \de\end{smallmatrix}\right] (\tau,0)=0$ with $\langle \ep , \de\rangle =0$ ${\rm mod}\, 2$. If none of the theta constants vanish, we can prove there always exist cubic relations that are not a product of a quadratic and a linear form; this is done by means of a dimensional argument, as we have
\[
 \dim_{\CC} S^3 B_1=2^g\big(2^g+1\big)\big(2^{g-1}+1\big)/3> 2^{g-1}\big( 3^g+1\big)=\dim_{\CC} S^3 B_3^+.
\]
 In the genus 3 case there are exactly 8 cubics: in the non-hyperelliptic case (no vanishing theta constants) these cubics can be obtained as derivatives of the Coble quartic, while in the hyperelliptic case (one vanishing theta constant) they are obtained as a product of a quadric $Q[\ep, \ep'](\tau, z)$ and a linear form in the $\T[\s](\tau,z)$.

When the genus is higher, we need to determine a set of theta constants $\theta \left[\begin{smallmatrix} \ep \\ \ep ' \end{smallmatrix}\right] (\tau,0)$ that vanish at the point $\tau$ in the hyperelliptic locus; to do this, we start with a hyperelliptic point to which we can associate a special fundamental system of characteristics, namely an azygetic collection of $2g+2$ characteristics $m_1, \dots, m_g, m_{g+1}, \dots, m_{2g+2}$ such that the first $g$ are odd and the last $g+2$ are even. Once we denote by $\{e_k\}_{k=1, \dots, g}$ the elements of the natural basis in ${\ZZ^g_2}$, and set $e_{g+1}= 0$ and $s_k= e_1+\dots+e_k$, we can choose the following as special fundamental system
\begin{gather}
 m_1:=\begin{bmatrix} s_1 \\ e_1\end{bmatrix},\qquad \dots, \qquad m_k:=\begin{bmatrix} s_k \\ e_k\end{bmatrix},\qquad\dots,\qquad m_g:=\begin{bmatrix} s_g \\ e_g\end{bmatrix},
\nonumber\\
m_{g+1}:=\begin{bmatrix} 0 \\ e_{1}\end{bmatrix},\qquad \dots,\qquad m_{2g}:=\begin{bmatrix} s_{g-1} \\ e_{g}\end{bmatrix}, \qquad m_{2g+1}:=\begin{bmatrix} s_g \\ 0
\end{bmatrix},\qquad m_{\infty}:=\begin{bmatrix} 0 \\ 0 \end{bmatrix} .\label{eq:FS}
\end{gather}

To be consistent with Mumford's notation in \cite[p.~106]{mu}, we set
\begin{gather}\label{BU}
B:=\{1, \dots, 2g+ 1, \infty \} ,\qquad U:= \{ g+1, \dots , 2g+1 \},
\end{gather}
and for any subset $S\subset B$ we denote by $ CS$ the complementary set in $B$.
 To an even subset $T\subset B$ (modulo $S \sim CS$) we associate a characteristic according to the rule
 \[
T\to m_T= \sum_{j\in T} m_j .
\]
This actually defines a bijection.

By setting $T\circ S:= (T\setminus S)\cup (S\setminus T)$, it can be checked that the parity of $m_T$ is equal to $(-1)^{\frac{\#(T\circ U)-g-1}{2}}$.

We then know from the characterization of the hyperelliptic locus (cf.~\cite{Th} and~\cite{mu}) that the period matrix of a hyperelliptic curve always admits a conjugate $\tau \in \HH_g$ under the action of the group ${\rm Sp}(2g, \ZZ)$, for which the following vanishing and non-vanishing conditions for theta constants hold
\begin{gather}
\theta_{m_T} (\tau,0)=0 \qquad \text{if} \quad \#(T\circ U)\neq g+1, \label{vanishing}\\
\theta_{m_T} (\tau,0) \neq 0 \qquad \text{if} \quad \#(T\circ U) = g+1. \label{nonvanishing}
\end{gather}
